% Appended to Part C to preserve previously assigned statement numbers.
\begin{proposition}[Proved: a finite rational primitive criterion]
\label{C4:rational-primitive-criterion}
Let $s\geq2$ and let $a_0,\ldots,a_k:P\to Q$ be a finite family
in $\mathcal A_s(\mathbb Q)$. The following are equivalent:
\begin{enumerate}
\item the family has an affine equalizer over $\mathbb Q$;
\item there are primitives $b_P\in C^{s-1}(\Delta^P;\mathbb Q)$
and $b_Q\in C^{s-1}(\Delta^Q;\mathbb Q)$ of the specified cocycles
such that $a_i^*b_Q=b_P$ for every $i$;
\item the family has an affine equalizer over $\mathbb Z$.
\end{enumerate}
Thus rational equalizability for a fixed family is a finite system
of rational linear equations on its two given objects. This
assertion concerns the coefficient group $\mathbb Q$ of the
cocycles; it is not asserted with a finite coefficient group in its
place.
\end{proposition}
\begin{proof}
Suppose $t=\sum_g\lambda_g g\in
\mathbb Q[\mathcal A_s(\mathbb Q)(Q,D)]$ is an affine equalizer.
Choose a primitive $b_D$ of the cocycle on the simplex $D$ and set
\[
 b_Q=\sum_g\lambda_g g^*b_D.
\]
Because $\sum_g\lambda_g=1$, its coboundary is the cocycle on
$Q$. The equality of the formal composites $ta_i$ makes all
$a_i^*b_Q$ equal; their common value is the required $b_P$.
Conversely such primitives lift the whole family to one pair of
objects in the primitive-cochain expansion. That expansion is the
Ramsey class of ordered arbitrary $(s-1)$-cochains used in
Theorem~\ref{C4:counterexample}. Its integral affine criterion gives
an integral equalizer of the lifted family, and forgetting the
primitives gives one in $\mathcal A_s(\mathbb Q)$. Finally an
integral equalizer is a rational one. The converse uses the fact
that the barycenters above still take values in the actual cocycle
coefficient group $\mathbb Q$.
\end{proof}

\begin{theorem}[Proved: increasing cocycle degrees erase fixed families]
\label{C4:degree-colimit-exclusion}
Let $r,s\geq2$ and let
$F:\mathcal A_r(\mathbb Q)\to\mathcal A_s(\mathbb Q)$ be any
functor. If $a_0,\ldots,a_k:P\to Q$ is a finite parallel family
and $s\geq |Q|$, the family $F(a_0),\ldots,F(a_k)$ has an integral
affine equalizer. In particular, for any sequence of functors
\[
 \mathcal A_{r_0}(\mathbb Q)\longrightarrow
 \mathcal A_{r_1}(\mathbb Q)\longrightarrow\cdots,
 \qquad r_n\longrightarrow\infty,
\]
its filtered colimit as a small category is integrally exact.
Consequently this construction cannot give a total all-prime-exact
category which fails rational exactness, regardless of the chosen
transition functors.
\end{theorem}
\begin{proof}
We work with the equivalent category of finite subsets of the
Fraisse limit $M$ of $\mathcal A_r(\mathbb Q)$ and write
$G=\operatorname{Aut}(M)$. The previously checked cocycle-class
properties give (Op) and
$\operatorname{Fix}_M(G_B)=B$ for every finite $B\subset M$.
All finite subsets are strong, and the points of $M$ form one
$G$-orbit.

Take the filtered colimit of the underlying ordered sets of $F(B)$
over inclusions of finite nonempty subsets $B\subset M$. It is a
nonempty ordered smooth $G$-set $S$, since every transition map is
an increasing injection. Its $G$-action is induced by
$F(g:B\to gB)$. The cocycles on the $F(B)$ agree under inclusions,
so give a $G$-invariant cocycle
\[
 c\in C^s(\Delta^S;\mathbb Q).
\]
Here a finite tuple of vertices of $S$ lies in some $F(B)$, so this
definition covers every simplex and its coboundary equation.

If $S$ has a $G$-fixed vertex, coning chains to that vertex gives
a $G$-invariant primitive of $c$ on the entire simplex. Its
restrictions to $F(P)$ and $F(Q)$ satisfy the preceding proposition,
and we are done. Suppose therefore that $S$ has no fixed vertex.
For each $v\in S$, (Op) writes its stabilizer uniquely as
$G_{B_v}$ for a finite nonempty subset $B_v\subset M$. Uniqueness
follows from the displayed fixed-point identity. Moreover
\[
 B_{gv}=gB_v,
 \qquad v\in F(B)\ \Longrightarrow\ B_v\subset B,
 \tag{DC.1}
\]
because $G_B\subset G_v=G_{B_v}$ in the second assertion.
Choose $x\in M$ and a vertex $v\in F(\{x\})$. By (DC.1),
$B_v=\{x\}$, so its orbit supplies a $G$-equivariant injection
$j:M\to S$ with $j(x)=v$.

On rational vertex spaces define the augmentation-preserving,
$G$-equivariant maps
\[
 p:\mathbb Q[S]\longrightarrow\mathbb Q[M],
 \quad p(v)=\frac1{|B_v|}\sum_{x\in B_v}x,
 \qquad h=jp:\mathbb Q[S]\longrightarrow\mathbb Q[S].
\]
Use oriented simplex chains, identifying
$C_n(\Delta^S;\mathbb Q)$ with
$\bigwedge^{n+1}\mathbb Q[S]$ and the boundary with contraction
by augmentation. On each increasing basis tuple put
\[
\begin{split}
 H_n(v_0\wedge\cdots\wedge v_n)
 =\sum_{i=0}^n(-1)^i\,&h(v_0)\wedge\cdots\wedge h(v_i)\wedge v_i
 \\[-2pt]
 &{}\wedge\cdots\wedge v_n .
\end{split}
\tag{DC.2}
\]
Extend this formula linearly on the increasing-tuple basis. All
terms have finite support. The ordinary prism telescoping calculation
gives
\[
 \partial H+H\partial=\operatorname{id}-\bigwedge h.
 \tag{DC.3}
\]
For example, in degree zero it is
$\partial(h(v)\wedge v)=v-h(v)$; in higher degrees the internal
faces cancel in consecutive summands. The same calculation is valid
for rational barycenters because every $h(v)$ has augmentation one.
The use of increasing tuples is essential to the definition:
$G$ preserves their order, and hence (DC.2) is $G$-equivariant.
No invariance under arbitrary reordered tuples is being presumed.

Define the $G$-invariant cochain $b=cH_{s-1}$. Equation (DC.3) and
$dc=0$ give
\[
 db=c-(\bigwedge h)^*c.
\]
For vertices in $F(Q)$, (DC.1) implies that their images under $h$
lie in the vector space spanned by $j(Q)$, which has dimension
$|Q|$. Thus the second term vanishes on degree-$s$ chains of
$F(Q)$ whenever $s\geq |Q|$. Consequently $b|_{F(Q)}$ is a
primitive of its cocycle. The same holds for $F(P)$, since
$|P|\leq |Q|$. Every $a_i$ extends to an element of $G$ by
homogeneity. The $G$-invariance of $b$ therefore gives
\[
 F(a_i)^*(b|_{F(Q)})=b|_{F(P)}
 \qquad(0\leq i\leq k).
\]
Proposition~\ref{C4:rational-primitive-criterion} now supplies the
integral affine equalizer.

Finally every finite parallel family in a filtered colimit of
categories can be represented at one finite stage: its finitely
many objects, arrows, and required equalities all occur at some
common stage. Represent a family at stage $n$ with target $Q$,
and choose $m\geq n$ with $r_m\geq |Q|$. Apply the first assertion
to the composite functor from stage $n$ to stage $m$ and then
pass its integral equalizer to the colimit. Common successors and
nonemptiness also pass to this colimit. The all-family affine
criterion proves integral exactness. This proof does not assume
that the transition functors are faithful or injective on objects.
\end{proof}

